% Appendix: Model Onboarding Under Budget Constraints
% Include from paper/sections/appendix.tex

\subsection{Model Onboarding Under Budget Constraints}
\label{appendix:model_onboarding}

Production LLM portfolios are not static.
Providers release new models monthly---Gemini~2.5 Flash launched five
months after Gemini~2.5 Pro, GPT-4o-mini arrived three months after
GPT-4o---and platform teams need to incorporate them without
disrupting cost compliance or quality.
This section evaluates whether a single API call,
\texttt{register\_model()}, suffices to onboard a new model into a
\emph{budget-constrained} router---and whether the router correctly
\emph{rejects} models that are inferior or too expensive.

\subsubsection{Why Care?}

The standard industry workflow for integrating a new model into a
supervised router is:
\begin{enumerate}[nosep,leftmargin=*]
  \item Collect the new model's responses on a representative prompt
    set and label them (days to weeks of human or LLM-judge
    annotation).
  \item Retrain the routing policy on the expanded labeled dataset.
  \item Update hard-coded cost tables and budget thresholds.
  \item Re-deploy and monitor for regressions.
\end{enumerate}
This process is slow, manual, and error-prone: stale cost tables
are the leading cause of budget violations in production
routers~(internal post-mortems).
Critically, step~1 requires a \emph{pre-collected evaluation dataset}
for the new model before it can receive any live traffic---an
offline data-collection campaign that must be repeated for every
model addition.

ParetoBandit replaces steps 1--4 with a single call that seeds the new
arm with a T-shirt prior and lets the online learner discover its
quality--cost profile from live traffic.
The bandit still requires a reward signal, but that signal comes
from the same evaluation mechanism (LLM judge, user ratings, or
implicit behavioural composites) already serving the existing
portfolio---no dedicated offline dataset is collected for the new
model.
When the exploration mechanism routes a prompt to the new arm, the
existing reward pipeline evaluates its response at zero marginal
data-collection cost, and the bandit updates its posterior
immediately.
This also avoids two risks inherent to offline evaluation:
\emph{distribution mismatch} (a curated eval set may not reflect the
live prompt distribution) and \emph{dataset staleness} (the eval ages
as deployment patterns drift).
The BudgetPacer automatically accommodates the new arm's cost
without any reconfiguration.

\subsubsection{Experimental Design}

The same K=3 portfolio, data splits, and warmup priors used in
Experiments~01--03 are deployed with the BudgetPacer.
After Phase~1 online learning on the validation split ($n{=}1{,}785$),
\texttt{register\_model()} adds Gemini-2.5-Flash as a fourth arm
with \texttt{speed="fast"} and its published pricing
(input: \$0.30/M, output: \$2.50/M, blended: \$1.40/M).
Phase~2 continues online on the K=4 test split ($n{=}1{,}793$),
where Flash rewards are judged by DeepSeek-R1 using the same v3
continuous rubric as all other arms.

\paragraph{Forced exploration (burn-in).}
With $\alpha{=}0.01$ and strong K=3 warmup priors
($n_\mathit{eff}{=}1{,}163.9$), the UCB exploration bonus for a
cold-start arm is too small to trigger natural exploration: Flash
starts with $A{=}I$, $b{=}0$ ($\theta{=}0$), and
$\alpha\sqrt{x^\top A^{-1} x}$ cannot compete with K=3 arms' learned
predictions (${\sim}0.5$--$0.7$).
To address this, the first 20 Phase~2 prompts are routed
unconditionally to Flash (forced exploration).
This guarantees the bandit collects real observations before UCB
selection takes over.
After burn-in, the algorithm has enough data to discriminate: a
good model's $\theta$ is competitive and UCB adopts; a bad model's
$\theta$ is low and UCB rejects.
Forced exploration is standard in the contextual bandit literature
(cf.\ ``forced exploration at decaying rate'' in LLM routing,
arXiv~2602.02061).

All conditions use the epsilon-constraint-tuned hyperparameters from
Experiment~05: $\alpha{=}0.01$,
$n_{\mathrm{eff}}{=}1{,}163.9$, $\gamma{=}0.997$.

\paragraph{Onboarding scenarios.}
Three scenarios test whether the router \emph{discriminates} rather
than blindly adopting every new model:
\begin{itemize}[nosep,leftmargin=*]
  \item \textbf{Good \& Cheap} (baseline): Flash as collected---high
    quality, low cost.  Expected: adopted.
  \item \textbf{Good \& Expensive}: Flash costs scaled $10\times$.
    Expected: suppressed under tight budgets by the pacer, partially
    adopted when unconstrained.
  \item \textbf{Bad \& Cheap}: Flash rewards scaled $0.5\times$.
    Expected: explored during burn-in, then rejected once the bandit
    learns the low reward.
\end{itemize}

\paragraph{Strategies.}
Two Phase~2 strategies are compared per budget tier:
\begin{itemize}[nosep,leftmargin=*]
  \item \textbf{Fixed Policy (uniform 1/4):}
    Equal allocation across all K=4 arms.
    The simplest possible onboarding strategy---no routing
    intelligence, no budget awareness.
  \item \textbf{ParetoBandit (transfer):}
    Phase~1 posteriors carry over; Flash gets a T-shirt prior
    plus 20-prompt forced burn-in.
    BudgetPacer continues without reconfiguration.
\end{itemize}

\noindent
Three budget targets span the constraint regime:
tight ($B{=}\$3.0{\times}10^{-4}$/req),
moderate ($B{=}\$6.6{\times}10^{-4}$/req),
and loose ($B{=}\$1.9{\times}10^{-3}$/req).
An unconstrained ParetoBandit baseline (no pacer) serves as a quality
ceiling.  Results: 20~seeds (9000--9019).

\subsubsection{Results (Figure~\ref{fig:model_onboarding})}

\paragraph{(a) Good \& Cheap: Flash is adopted across all budget tiers.}

\begin{center}
\small
\begin{tabular}{@{}llrrr@{}}
\toprule
\textbf{Strategy} & \textbf{Budget} &
  \textbf{Reward} & \textbf{Cost (\$/req)} &
  \textbf{Flash \%} \\
\midrule
Fixed Policy & any
  & 0.894 & \$4.06\text{e-}3 & 25.1\% \\
\midrule
ParetoBandit & tight
  & 0.866 & \$3.00\text{e-}4 & 8.6\% \\
ParetoBandit & moderate
  & 0.899 & \$6.52\text{e-}4 & 12.9\% \\
ParetoBandit & loose
  & 0.918 & \$1.69\text{e-}3 & 15.9\% \\
ParetoBandit & unconstrained
  & 0.929 & \$1.05\text{e-}2 & 13.3\% \\
\bottomrule
\end{tabular}
\end{center}

\noindent
Flash achieves sustained adoption ($>10$\% windowed share for 50
consecutive post-washout steps) in 20/20~seeds across all budget
tiers.
Under the loose budget, Flash stabilises at $10.2$\% of late-phase
traffic; even under the tight budget, it maintains $4.4$\%.
The Fixed Policy spends \$4.06e-3 per request---exceeding the tight
budget by $13.5\times$ and the loose budget by $2.2\times$---while
ParetoBandit stays within target at every tier.

\paragraph{(b) Good \& Expensive: the BudgetPacer correctly suppresses
costly models.}
When Flash costs are scaled $10\times$ (\$3.00/M input,
\$25.00/M output), the pacer's hard ceiling filter and soft cost
penalty correctly identify it as expensive.
Under tight and moderate budgets, Flash never achieves sustained
adoption (0/20~seeds).
Under loose budgets, Flash partially recovers ($6.1$\% final share,
sustained in 20/20~seeds), and when unconstrained, it reaches
$7.3$\%---demonstrating that the pacer, not the bandit's quality
estimate, drives the suppression.

\paragraph{(c) Bad \& Cheap: the bandit correctly rejects low-quality models.}
With rewards halved, Flash's $\theta$ after burn-in is
uncompetitive.
In every seed and every budget tier, Flash's final share drops
to exactly $0.0$\%: the 20-prompt burn-in costs
${\sim}1$\% of Phase~2 steps, after which UCB never selects Flash
again.
This confirms that forced exploration is self-limiting: a bad model
incurs bounded $O(N)$ regret during burn-in and is permanently
rejected once the posterior stabilises.

\paragraph{Fixed Policy: budget-oblivious baseline.}
The uniform-1/4 baseline highlights the cost of ignoring budget
constraints.
In the Good \& Cheap scenario, it spends \$4.06e-3 per
request---exceeding the tight budget by $13.5\times$.
In the Bad \& Cheap scenario, giving 25\% of traffic to a
half-quality model drags overall reward to 0.779, compared to
ParetoBandit's 0.925 (unconstrained).
In the Good \& Expensive scenario, the Fixed Policy's cost
balloons to \$6.07e-3 per request.

\subsubsection{Practical Implications}

\paragraph{One API call, zero downtime.}
At moderate and loose budgets---the regime of most production
deployments---onboarding a new model into ParetoBandit requires a
single call to \texttt{register\_model()} with the model's pricing,
plus a brief forced-exploration period (20~prompts, ${\sim}1$\% of
traffic).
No offline evaluation, no policy retraining, no budget
reconfiguration.  The bandit discovers the model's quality--cost
niche from live traffic and the pacer auto-adjusts.

\paragraph{Discrimination, not blind adoption.}
The three-scenario design demonstrates that ParetoBandit does not
simply absorb every new model.
Good models are adopted, expensive models are budget-gated, and bad
models are rejected---all from the same \texttt{register\_model()}
call with no scenario-specific tuning.

\paragraph{Tight budgets require deliberate exploration.}
Under tight budget constraints, the pacer's exploration budget is
limited, slowing Flash adoption.
The Good \& Expensive scenario shows that tight budgets
can effectively block an expensive newcomer (0/20~seeds sustain),
while loose budgets permit partial adoption.
Practitioners running at very tight margins who want to onboard an
expensive newcomer should temporarily increase the budget target to
allow initial exploration, then tighten once the bandit calibrates.
